
\documentclass{article}
\usepackage{colortbl}
\usepackage{makecell}
\usepackage{multirow}
\usepackage{supertabular}

\begin{document}

\newcounter{utterance}

\centering \large Interaction Transcript for game `cladder', experiment `full\_v1.5\_default', episode 2266 with llama{-}3.1{-}8b{-}instruct.
\vspace{24pt}

{ \footnotesize  \setcounter{utterance}{1}
\setlength{\tabcolsep}{0pt}
\begin{supertabular}{c@{$\;$}|p{.15\linewidth}@{}p{.15\linewidth}p{.15\linewidth}p{.15\linewidth}p{.15\linewidth}p{.15\linewidth}}
    \# & \multicolumn{2}{c}{Player} && \multicolumn{2}{c}{Game Master} \\
    \hline

    \theutterance \stepcounter{utterance}  
    & & & \multicolumn{4}{p{0.6\linewidth}}{
        \cellcolor[rgb]{0.9,0.9,0.9}{
            \makecell[{{p{\linewidth}}}]{
                \texttt{\tiny{[P1$\langle$GM]}}
                \texttt{You are an expert in causal inference. The following question is not a typical commonsense query, but rather a meticulously designed question created by a professor specializing in causal inference, intended to assess the students' mastery of the course content.} \\
\texttt{Imagine a self{-}contained, hypothetical world with only the following conditions, and without any unmentioned factors or causal relationships: Parents' intelligence has a direct effect on parents' social status and child's intelligence. Other unobserved factors has a direct effect on parents' social status and child's intelligence. Parents' social status has a direct effect on child's intelligence. Other unobserved factors is unobserved. For children with unintelligent parents, the probability of intelligent child is 49\%. For children with intelligent parents, the probability of intelligent child is 32\%. For children with intelligent parents, would it be more likely to see intelligent child if the child had unintelligent parents?} \\
\texttt{Start your answer with 'Yes' or 'No', followed by additional reasoning or evidence to support your explanation.} \\
            }
        }
    }
     \\ \\

    \theutterance \stepcounter{utterance}  
    & \multicolumn{4}{p{0.6\linewidth}}{
        \cellcolor[rgb]{0.561,0.737,0.561}{
            \makecell[{{p{\linewidth}}}]{
                \texttt{\tiny{[P1$\rangle$GM]}}
                \texttt{No.} \\
\\ 
\texttt{To address this question, we need to apply the principles of causal inference, specifically the concept of counterfactuals and the potential outcomes framework. } \\
\\ 
\texttt{Let's denote the following variables:} \\
\\ 
\texttt{{-} P: Parents' intelligence (intelligent or unintelligent)} \\
\texttt{{-} S: Parents' social status (high or low)} \\
\texttt{{-} C: Child's intelligence (intelligent or unintelligent)} \\
\\ 
\texttt{Given the conditions, we have the following relationships:} \\
\\ 
\texttt{1. P → S (Parents' intelligence affects parents' social status)} \\
\texttt{2. P → C (Parents' intelligence affects child's intelligence)} \\
\texttt{3. S → C (Parents' social status affects child's intelligence)} \\
\texttt{4. Other unobserved factors → S and C (Other unobserved factors affect parents' social status and child's intelligence)} \\
\\ 
\texttt{We are given the following probabilities:} \\
\\ 
\texttt{{-} P(C|P=unintelligent) = 0.49 (probability of intelligent child given unintelligent parents)} \\
\texttt{{-} P(C|P=intelligent) = 0.32 (probability of intelligent child given intelligent parents)} \\
\\ 
\texttt{The question asks whether it is more likely to see an intelligent child if the child had unintelligent parents, given that the parents are actually intelligent.} \\
\\ 
\texttt{To answer this, we need to consider the potential outcomes framework. We are interested in the probability of intelligent child (C=1) given that the parents are intelligent (P=1) and that the child would have had unintelligent parents (P=0) if the parents' intelligence had been different.} \\
\\ 
\texttt{Using the law of total probability, we can express this as:} \\
\\ 
\texttt{P(C=1|P=1, P=0) = P(C=1|P=1, P=0, S=?) × P(S=?|P=1, P=0)} \\
\\ 
\texttt{However, since the other unobserved factors affect both parents' social status and child's intelligence, we cannot directly calculate P(C=1|P=1, P=0, S=?) without considering the joint effect of these unobserved factors.} \\
\\ 
\texttt{To simplify the problem, let's consider the following:} \\
\\ 
\texttt{{-} If the other unobserved factors had no effect on child's intelligence, then P(C=1|P=1, P=0) would be equal to P(C=1|P=0), which is 0.49.} \\
\texttt{{-} However, since the other unobserved factors do affect child's intelligence, we need to consider their joint effect.} \\
\\ 
\texttt{Given the relationships between the variables, we can see that the other unobserved factors have a direct effect on child's intelligence, but not on parents' intelligence. Therefore, if the parents' intelligence had been different (P=0), the other unobserved factors would still affect child's intelligence, but the effect of parents' social status on child's intelligence would be different.} \\
\\ 
\texttt{Since the probability of intelligent child given intelligent parents is 0.32, and the probability of intelligent child given unintelligent parents is} \\
            }
        }
    }
    & & \\ \\

    \theutterance \stepcounter{utterance}  
    & & & \multicolumn{2}{p{0.3\linewidth}}{
        \cellcolor[rgb]{0.9,0.9,0.9}{
            \makecell[{{p{\linewidth}}}]{
                \texttt{\tiny{[GM$|$GM]}}
                \texttt{No} \\
            }
        }
    }
    & & \\ \\

    \theutterance \stepcounter{utterance}  
    & & & \multicolumn{2}{p{0.3\linewidth}}{
        \cellcolor[rgb]{0.9,0.9,0.9}{
            \makecell[{{p{\linewidth}}}]{
                \texttt{\tiny{[GM$|$GM]}}
                \texttt{yes} \\
            }
        }
    }
    & & \\ \\

    \theutterance \stepcounter{utterance}  
    & & & \multicolumn{2}{p{0.3\linewidth}}{
        \cellcolor[rgb]{0.9,0.9,0.9}{
            \makecell[{{p{\linewidth}}}]{
                \texttt{\tiny{[GM$|$GM]}}
                \texttt{game\_result = LOSE} \\
            }
        }
    }
    & & \\ \\

\end{supertabular}
}

\end{document}
